\documentclass[10pt,a4paper]{amsart}
\usepackage[margin=19mm]{geometry}
\usepackage{amsmath,amssymb,mathtools}
\usepackage[scr=rsfs,cal=euler]{mathalfa}
\usepackage{xfrac}
\usepackage[unicode,hidelinks,bookmarks=false]{hyperref}
\newcommand{\ie}{i.e.\ }
\newcommand{\cf}{cf.\ }
\newcommand{\resp}{respectively\ }
\newcommand{\Subsection}{\subsection}
\providecommand{\nobreakdash}{\nobreak\hbox{-}}
\setlength{\emergencystretch}{2em}
\multlinegap0pt
\usepackage{bm}
\usepackage{bbm}
\usepackage{dsfont}
\usepackage{xspace}
\usepackage{cancel}
\usepackage{mathtools}
\usepackage{amsthm}
\makeatletter
\@ifundefined{ctheorem}{\newtheorem{ctheorem}{Ctheorem}}{}
\@ifundefined{defn}{\newtheorem{defn}[ctheorem]{Definition}}{}
\@ifundefined{restatable}{\newtheorem{restatable}{Restatable}}{}
\makeatother
\DeclareMathOperator{\twostep}{2-step}
\title[Families of tractable problems with respect to vertex-interv]{Families of tractable problems with respect to vertex-interval-membership width and its generalisations}
\author{Jessica Enright, Samuel D. Hand, Laura Larios-Jones, Kitty Meeks}
\date{Source: arXiv:2505.15699v5 (CC BY 4.0)}
\begin{document}
\maketitle

\noindent\emph{Source and licence.} Families of tractable problems with respect to vertex-interval-membership width and its generalisations. Version arXiv:2505.15699v5 (CC BY 4.0), distributed under the Creative Commons Attribution 4.0 International licence, as recorded in the arXiv licence metadata for this exact version and shown on the abstract page https://arxiv.org/abs/2505.15699v5. Full rights notice: licenses/arxiv-2505.15699-CC-BY-4.0.txt.\par\smallskip

\noindent\textit{Editorial scope.} Counted unit: the Lemma whose source label is \texttt{TIMrootrealisable} in the article (source0.tex lines 1541--1543, source label \texttt{TIMrootrealisable}), delivered together with the article's own complete original proof of it (source0.tex lines 1544--1563, one \texttt{proof} environment of 20 lines). No mathematics is written, repaired or re-derived: statement and proof are the article's own text line for line. Required context retained uncounted: def:realisable (lines 1475--1488). Every definition, display and label of those lines is retained unchanged. Omitted from this package and not redistributed: 1--1474, 1489--1540, 1564--2537. Nothing inside a retained excerpt is omitted or altered: every retained body is delivered exactly as the article's lines are written, so no omission receipt of any kind accompanies this package. Citations: the retained excerpts carry no citation command at all, so no bibliography entry of the article is transcribed and no reference is redistributed. Reference adaptation: none was needed and none was made; every reference inside the delivered unit resolves inside it, so no unresolved reference or citation is produced. Printed slips kept literal: no printed word, symbol, hypothesis or number of the counted statement or of its proof was altered. Layout and class: the article's own class is \texttt{lipics-v2021} with packages including inputenc, amsmath, amsfonts, amsthm, amssymb, framed, xcolor, hyperref, complexity, xspace, booktabs, algorithm, algpseudocodex, bold-extra; the wrapper is the lane's pinned \texttt{amsart} editorial wrapper at 10pt on A4, which restates the theorem-like environments the retained text needs and adds the article's own macro definitions that the retained text needs (\texttt{\textbackslash twostep}), with extra wrapper packages (bm, bbm, dsfont, xspace, cancel, mathtools). The delivered numbering is the unit's own. Assets: no retained span contains a figure environment, a graphics-inclusion command, a PDF-page inclusion, a table or a TikZ picture; no image file is copied into this package, the package declares no asset dependency and no image file is redistributed with it. Licence: version arXiv:2505.15699v5 of this article is distributed under the Creative Commons Attribution 4.0 International licence, as recorded in the arXiv licence metadata for this exact version and shown on the abstract page https://arxiv.org/abs/2505.15699v5. A full rights notice is delivered at licenses/arxiv-2505.15699-CC-BY-4.0.txt. Characters: every retained excerpt is pure ASCII, so the delivered engine, XeLaTeX with its default Latin Modern text font, reports no missing character.
\begin{defn}\label{def:realisable}
    A $\twostep$ $(X,k)$ profile $\sigma=(l^{t-1}, l^t, l^{t+1}, \textbf{v}_{1},\ldots,\textbf{v}_{|\mathcal{C}^s|}, \textbf{total})$ of a bag $B^2(s)$ is \emph{realisable} if and only if there exists a configuration $\varsigma$ which maps every timed connected component $(C,t)\in \mathfrak{C}^s$ to a $(X,k)$-state $\varsigma(C,t)=(l_{\varsigma}(C,t),\textbf{v}_{\varsigma}(C,t))$
    % : \mathcal{C}^s\to \mathfrak{L}_X\times \mathbb{Z}^k$ of $\mathcal{G}^s$, where $\varsigma(C)=(\mathfrak{l}^t_C,\mathfrak{v}^t_C)$, $\mathfrak{L}_X$ is the set of all labellings $\mathfrak{l}^t_C$ which label the vertex-time pairs of the connected component $C$ at time $t$ with elements in $X$
    such that
    \begin{itemize}
        \item for every connected component-time pair $(C,t)$ in $B^2(s)$, $\sigma|_C=\varsigma(C,t)$,
        \item for all connected components $C$ of $G_1^s$, $\textbf{St}(\varsigma(C,0),C,x)=\textbf{true}$,
        \item for all connected components $C$ of $G^s_{\Lambda(\mathcal{G})}$, $\textbf{Fin}(\varsigma(C,\Lambda(\mathcal{G})),C,x)=\textbf{true}$,
        \item for all times $t\in(0,\Lambda(\mathcal{G}))$, every connected component $C$ of $G_{t}^s$, $\textbf{Val}(\varsigma(C,t),C,x)=\textbf{true}$,
        \item for all pairs $(C,t)$ in $\mathfrak{C}^s$, $\textbf{Tr}(l^{\varsigma}_{t-1}|_C,l_{\varsigma}(C,t),C,x)=\textbf{true}$ where $l^{\varsigma}_{t-1}$ is the labelling of all vertices in $G^s_{t-1}$ such that its restriction to any connected component $C\in G_{t-1}^s$ is $l_{\varsigma}(C,t-1)$, and
        \item $\textbf{total}= \sum_{0\leq t\leq \Lambda(\mathcal{G}^s)} \sum _{C \in \mathcal{C}(G^s_t)} \textbf{v}_{\varsigma}(C,t)$.
    \end{itemize}
    Here we say that $\varsigma$ \emph{realises} $\sigma$.
\end{defn}

\begin{restatable}{lemma}{TIMrootrealisable}\label{lem:TIM-root-realisable}
    Let $P$ be a $(X,k,f)$-component-exchangeable temporally uniform problem. Then an instance $x=(\mathcal{G},\beta)$ is a yes-instance of $P$ if and only if there is a realisable $\twostep$ $(X,k)$ profile $\sigma=(l^{t-1},l^t,l^{t+1},\textbf{v}_1,\ldots\textbf{v}_{|\mathcal{C}^r|}, \textbf{total})$ of the root bag $B^2(r)$ of the $\twostep$ TIM decomposition of $\mathcal{G}$ such that $\textbf{total}\leq \textbf{v}_{\text{upper}}$.
\end{restatable}

\begin{proof}
    We begin by supposing that $x=(\mathcal{G},\beta)$ is a yes-instance of $P$. Then, by definition of being $(X,k,f)$-component-exchangeable temporally uniform, there exists a sequence of $(X,k)$-component states $s_0,\ldots,s_{\Lambda}$ of the form $s_t=(l_t,\textbf{w}^t_1,\ldots,\textbf{w}^t_c, \nu_t)$ of each snapshot of $\mathcal{G}$ such that 
    \begin{itemize}
        \item for each connected component $C_1$ of $G_1$, $\textbf{St}(s_0|_{C_1},C_1,x)=\textbf{true}$;
        \item for each connected component $C_{\Lambda}$ of $G_{\Lambda}$, $\textbf{Fin}(s_{\Lambda}|_{C_{\Lambda}},C_{\Lambda},x)=\textbf{true}$;
        \item $\textbf{Tr}(l_{t-1}|_{C_t}, l_t|_{C_t}, C_t,x)=\textbf{true}$ where $l_t$ is the labelling of vertices of state $s_t$, for all times $1 \leq t \leq \Lambda$ and connected components $C_t$ of $G_t$;
        \item $\textbf{Val}(s_t|_{C_t}, C_t,x)=\textbf{true}$ for all times $1 < t < \Lambda$ and connected components $C_t$ of $G_t$; and
        \item the sum $\sum_{0\leq t\leq \Lambda} \sum_{C \in\mathcal{C}_t}\nu_{t} (C)\leq \textbf{v}_{\text{upper}}$.
    \end{itemize}
    From this, using Definition~\ref{def:realisable}, we can construct a configuration $\varsigma$ of $\mathcal{G}$ by letting $\varsigma(C,t)=s_t|_C$ for all connected components $C$ of all snapshots $G_t$ of $\mathcal{G}$. Then $\varsigma$ realises the profile $\sigma$ of $B^2(r)$ consisting of the restrictions of $l_t$ to the set of vertices $\{v \,:\,(v,t)\in B^2(s)\}$ for each time in the bag, the vectors assigned to each connected component in $B^2(r)$ by $\varsigma$, and the vector \textbf{total} which is given by the sum $\sum_{t\in [0,\Lambda]}\sum_{C\in\mathcal{C}(G_t)}\textbf{v}_{\varsigma}(C,t)$, where $\textbf{v}_{\varsigma}(C,t)$ is the vector given by $\varsigma(C,t)$.

    Now, suppose that there exists a $\twostep$ $(X,k)$ profile $\sigma=(l^{t-1},l^t,l^{t+1},\textbf{v}_1,\ldots\textbf{v}_{|\mathcal{C}^r|}, \textbf{total})$ of $B^2(r)$ for an instance $x$ of $P$ such that $\textbf{total}\leq \textbf{v}_{\text{upper}}$ and $\sigma$ is realisable. Let $\varsigma$ be the configuration of $\mathcal{G}$ which realises $\sigma$. Then we construct a sequence $s_0,\ldots, s_{\Lambda}$ of $(X,k)$-component states from $\varsigma$ by letting $l_{t}$ be the labelling that extends $l_{\varsigma}(C,t)$ for all connected components $C$ in $G_t$, and $\nu_t(C)=\textbf{v}_{\varsigma}(C,t)$, where $\textbf{v}_{\varsigma}(C,t)$ is the vector given by $\varsigma(C,t)$. Since the vector $\textbf{total}$ in $\sigma$ is the sum of all vectors of all timed connected components in $\mathcal{G}$ under $\varsigma$, $\textbf{total}$ must also be $\sum_{t\in[0,\Lambda]}\sum_{C\in\mathcal{C}_t}\nu_t(C)\leq \textbf{v}_{\text{upper}}$. Note that, for the root bag $B^2(r)$, $\mathfrak{C}^r$ is the set of all connected components in the graph ($\mathfrak{C}^r=\mathcal{C}(\mathcal{G}))$. Therefore, by definition, we must have that 
    \begin{itemize}
        \item for each connected component $C_1$ of $G_1$, $\textbf{St}(s_0|_{C_1},C_1,x)=\textbf{true}$;
        \item for each connected component $C_{\Lambda}$ of $G_{\Lambda}$, $\textbf{Fin}(s_{\Lambda}|_{C_{\Lambda}},C_{\Lambda},x)=\textbf{true}$;
        \item $\textbf{Tr}(l_{t-1}|_{C_t}, l_t|_{C_t}, C_t,x)=\textbf{true}$ where $l_t$ is the labelling of vertices of state $s_t$, for all times $1 \leq t \leq \Lambda$ and connected components $C_t$ of $G_t$; and
        \item $\textbf{Val}(s_t|_{C_t}, C_t,x)=\textbf{true}$ for all times $1 < t < \Lambda$ and connected components $C_t$ of $G_t$.
    \end{itemize}
    Therefore, $x$ must be a yes-instance of $P$. In conclusion, an instance $x$ is a yes-instance of a $(X,k,f)$-component-exchangeable temporally uniform problem if and only if there exists a realisable $\twostep$ $(X,k)$ profile of the root bag $B^2(r)$ of the $\twostep$ TIM decomposition of $\mathcal{G}$ such that $\textbf{total}\leq \textbf{v}_{\text{upper}}$.
\end{proof}

\end{document}
